% GENERATED FILE. Edit canonical lessons, then run npm run generate:books.
% canonical-source-hash: fnv1a64:4d787c735fb197bd
% canonical-lessons: SA-C179-sainikah, SA-C179-raksakah, SA-C179-nyayadhisah, SA-C179-adhivakta, SA-C179-lekhakah

\chapter{Soldier, Guard, Judge, Lawyer, Writer}
\label{ch:sa-soldier-guard-judge-lawyer-writer}
\hlchaptermodality{179}

\begin{chapteropening}
\textbf{By the end of this chapter:} \emph{I can say a soldier, a guard, a judge, a lawyer and a writer.}

Complete the last lesson of chapter 179: I can say a soldier, a guard, a judge, a lawyer and a writer.
\end{chapteropening}

\section[\texorpdfstring{\sk{सैनिकः} (sainikaḥ)}{सैनिकः (sainikaḥ)}]{\sk{सैनिकः} (sainikaḥ) — a soldier}

\label{lesson:SA-C179-sainikah}



\begin{quote}
Before the new one: say the Sanskrit for a son-in-law, then the Sanskrit for a daughter-in-law.
\end{quote}

\subsection*{You'll want to know: \sk{सैनिकः}}
\textbf{\sk{सैनिकः}} — \emph{sainikaḥ} — \textquotedblleft{}a soldier\textquotedblright{}.

\begin{grammarlens}[title={What you've built}]
One more word for this chapter.
\end{grammarlens}

\subsection*{Your turn}
\begin{itemize}
\raggedright
  \item \emph{Say it:} \emph{sainikaḥ}
  \item \emph{Say it:} \emph{sainikaḥ}, once more
  \item \emph{Say it:} say \emph{sainikaḥ}
  \item \emph{Recall:} say the Sanskrit for a son-in-law, then the Sanskrit for a daughter-in-law, then say \emph{sainikaḥ} again
\end{itemize}

\begin{tcolorbox}[breakable,colback=teal!4,colframe=teal!35!black,title={Before you move on}]
What does \sk{सैनिकः} mean? (\textquotedblleft{}A soldier\textquotedblright{}.) Say it once more.
\end{tcolorbox}
\section[\texorpdfstring{\sk{रक्षकः} (rakṣakaḥ)}{रक्षकः (rakṣakaḥ)}]{\sk{रक्षकः} (rakṣakaḥ) — a guard}

\label{lesson:SA-C179-raksakah}



\begin{quote}
Before the new one: say the Sanskrit for a daughter-in-law, then the Sanskrit for a soldier.
\end{quote}

\subsection*{You'll want to know: \sk{रक्षकः}}
\textbf{\sk{रक्षकः}} — \emph{rakṣakaḥ} — \textquotedblleft{}a guard\textquotedblright{}.

\begin{grammarlens}[title={What you've built}]
One more word for this chapter.
\end{grammarlens}

\subsection*{Your turn}
\begin{itemize}
\raggedright
  \item \emph{Say it:} \emph{rakṣakaḥ}
  \item \emph{Say it:} \emph{rakṣakaḥ}, once more
  \item \emph{Say it:} say \emph{rakṣakaḥ}
  \item \emph{Recall:} say the Sanskrit for a daughter-in-law, then the Sanskrit for a soldier, then say \emph{rakṣakaḥ} again
\end{itemize}

\begin{tcolorbox}[breakable,colback=teal!4,colframe=teal!35!black,title={Before you move on}]
What does \sk{रक्षकः} mean? (\textquotedblleft{}A guard\textquotedblright{}.) Say it once more.
\end{tcolorbox}
\section[\texorpdfstring{\sk{न्यायाधीशः} (nyāyādhīśaḥ)}{न्यायाधीशः (nyāyādhīśaḥ)}]{\sk{न्यायाधीशः} (nyāyādhīśaḥ) — a judge}

\label{lesson:SA-C179-nyayadhisah}



\begin{quote}
Before the new one: say the Sanskrit for a soldier, then the Sanskrit for a guard.
\end{quote}

\subsection*{You'll want to know: \sk{न्यायाधीशः}}
\textbf{\sk{न्यायाधीशः}} — \emph{nyāyādhīśaḥ} — \textquotedblleft{}a judge\textquotedblright{}.

\begin{grammarlens}[title={What you've built}]
One more word for this chapter.
\end{grammarlens}

\subsection*{Your turn}
\begin{itemize}
\raggedright
  \item \emph{Say it:} \emph{nyāyādhīśaḥ}
  \item \emph{Say it:} \emph{nyāyādhīśaḥ}, once more
  \item \emph{Say it:} say \emph{nyāyādhīśaḥ}
  \item \emph{Recall:} say the Sanskrit for a soldier, then the Sanskrit for a guard, then say \emph{nyāyādhīśaḥ} again
\end{itemize}

\begin{tcolorbox}[breakable,colback=teal!4,colframe=teal!35!black,title={Before you move on}]
What does \sk{न्यायाधीशः} mean? (\textquotedblleft{}A judge\textquotedblright{}.) Say it once more.
\end{tcolorbox}
\section[\texorpdfstring{\sk{अधिवक्ता} (adhivaktā)}{अधिवक्ता (adhivaktā)}]{\sk{अधिवक्ता} (adhivaktā) — a lawyer}

\label{lesson:SA-C179-adhivakta}



\begin{quote}
Before the new one: say the Sanskrit for a guard, then the Sanskrit for a judge.
\end{quote}

\subsection*{You'll want to know: \sk{अधिवक्ता}}
\textbf{\sk{अधिवक्ता}} — \emph{adhivaktā} — \textquotedblleft{}a lawyer\textquotedblright{}.

\begin{grammarlens}[title={What you've built}]
One more word for this chapter.
\end{grammarlens}

\subsection*{Your turn}
\begin{itemize}
\raggedright
  \item \emph{Say it:} \emph{adhivaktā}
  \item \emph{Say it:} \emph{adhivaktā}, once more
  \item \emph{Say it:} say \emph{adhivaktā}
  \item \emph{Recall:} say the Sanskrit for a guard, then the Sanskrit for a judge, then say \emph{adhivaktā} again
\end{itemize}

\begin{tcolorbox}[breakable,colback=teal!4,colframe=teal!35!black,title={Before you move on}]
What does \sk{अधिवक्ता} mean? (\textquotedblleft{}A lawyer\textquotedblright{}.) Say it once more.
\end{tcolorbox}
\section[\texorpdfstring{\sk{लेखकः} (lekhakaḥ)}{लेखकः (lekhakaḥ)}]{\sk{लेखकः} (lekhakaḥ) — a writer}

\label{lesson:SA-C179-lekhakah}



\begin{quote}
Before the new one: say the Sanskrit for a judge, then the Sanskrit for a lawyer.
\end{quote}

\subsection*{You'll want to know: \sk{लेखकः}}
\textbf{\sk{लेखकः}} — \emph{lekhakaḥ} — \textquotedblleft{}a writer\textquotedblright{}.

\begin{grammarlens}[title={What you've built}]
One more word for this chapter.
\end{grammarlens}

\subsection*{Your turn}
\begin{itemize}
\raggedright
  \item \emph{Say it:} \emph{lekhakaḥ}
  \item \emph{Say it:} \emph{lekhakaḥ}, once more
  \item \emph{Say it:} say \emph{lekhakaḥ}
  \item \emph{Recall:} say the Sanskrit for a judge, then the Sanskrit for a lawyer, then say \emph{lekhakaḥ} again
\end{itemize}

\begin{tcolorbox}[breakable,colback=teal!4,colframe=teal!35!black,title={Before you move on}]
What does \sk{लेखकः} mean? (\textquotedblleft{}A writer\textquotedblright{}.) Say it once more.
\end{tcolorbox}
